\section{Dispatch 层：\texttt{vma.c} 实现}

\texttt{vma.c} 的职责非常简单——校验后端是否已注册，然后转发：

\codefile{src/kernel/mm/vma.c}
\begin{lstlisting}[style=cstyle]
static const vma_ops_t* g_vma_ops = NULL;
static int g_vma_ready = 0;

void vma_register_backend(const vma_ops_t* ops) {
    g_vma_ops = ops;
}

void vma_init(void) {
    if (!g_vma_ops) {
        printk("[vma] no backend registered\n");
        return;
    }
    g_vma_ops->init();
    g_vma_ready = 1;
    printk("[vma] backend: %s, init ok\n", g_vma_ops->name);
}

int vma_add(uint32_t start, uint32_t end,
            uint32_t flags, const char* name) {
    if (!g_vma_ops || !g_vma_ready) return -1;
    return g_vma_ops->add(start, end, flags, name);
}
\end{lstlisting}

其余 \texttt{vma\_remove}/\texttt{vma\_find}/\texttt{vma\_dump}/\texttt{vma\_count}
格式完全一致—— \texttt{if (!ready) return fallback; return g\_vma\_ops->xxx(...);}。

\begin{designbox}
Dispatch 层看似"什么都没做"，但它带来三个好处：
\begin{enumerate}
  \item \textbf{编译期解耦}——调用者只 \texttt{\#include "vma.h"}，不需要知道后端头文件；
  \item \textbf{运行时切换}——\texttt{vma\_register\_backend} 在 init 前随时可换后端；
  \item \textbf{安全守卫}——未初始化时所有 API 返回安全默认值，不会空指针崩溃。
\end{enumerate}
\end{designbox}

% ============================================================================

